The increasing integration of automated employment decision tools (AEDTs) into hiring processes has significantly transformed how job applicants are evaluated, particularly in technology-driven fields. These tools—ranging from resume screening algorithms to automated coding assessments and AI-driven interviews—are often positioned as efficient and unbiased alternatives to traditional hiring practices. However, despite their growing adoption, there remains limited understanding of how applicants perceive these systems and adapt their strategies to navigate them. This gap is critical, as applicants are not merely passive participants but active agents who develop tactics such as resume optimization, coding practice, and leveraging referrals to improve their chances of success.

This report examines the intersection of applicant perceptions, strategy use, and hiring outcomes in the context of automated hiring systems. Prior research has highlighted concerns around transparency, fairness, and accountability in AEDTs, including evidence that such systems can perpetuate biases similar to human decision-makers. Moreover, applicants often experience distrust and a sense of dehumanization when evaluated algorithmically, particularly when processes lack clarity or human oversight. By analyzing how young job seekers perceive fairness, how they strategically respond to automation, and how these factors relate to job success, this work aims to provide a more nuanced understanding of the real-world implications of automated hiring and inform the design of more equitable and transparent systems.

\section{Motivation}
The use of automated employment decision tools (AEDTs) has rapidly increased in hiring contexts, especially for computing jobs. This includes the use of software in reading resumes, scoring coding assessments, and even conducting interviews to assess job applicants among others \cite{Bogen2018HelpWA,sanchez2020mean}. The motivation for this work is rooted in the rapid transformation of hiring practices through the adoption of automated employment decision tools (AEDTs), particularly in computing-related roles where large applicant pools necessitate scalable evaluation methods. These systems are now routinely used for tasks such as resume parsing, coding assessment evaluation, and even automated interviews \cite{sanchez2020mean}. While AEDTs are often promoted as efficient and unbiased, prior research has raised significant concerns about their \textbf{lack of transparency, accountability, and auditing}, as well as the limited visibility applicants have into how decisions are made \cite{3_ajunwa2024auditing,62_sloane2022silicon}. At the same time, regulatory efforts such as auditing requirements for AEDTs signal growing recognition of these issues, but there remains insufficient empirical understanding of how these systems are experienced by job seekers themselves. This creates a critical gap between system deployment and user-centered evaluation, motivating the need to examine applicant perspectives more closely.

Another key driver of this work is the evolving role of applicants as active participants in automated hiring ecosystems. Rather than passively undergoing evaluation, job seekers develop strategies to adapt to these systems, such as using resume scanners, practicing coding problems, or leveraging referrals to bypass automated filters \cite{4_armstrong2023navigating}. However, prior qualitative research suggests that applicants often feel constrained in their ability to fully demonstrate their skills in automated settings, especially compared to human evaluators who can interpret nuance and context \cite{4_armstrong2023navigating}. Furthermore, there exists a broader disconnect between employer priorities—such as efficiency and standardization—and applicant expectations for recognition, fairness, and meaningful evaluation \cite{26_giannakos2017understanding}. This tension highlights the need to systematically investigate how different strategies interact with perceptions of fairness and ultimately influence hiring outcomes.

Equally important is the growing body of evidence indicating that \textbf{algorithmic decision-making is frequently perceived as less fair than human judgment}, which can undermine trust in automated hiring systems and discourage applicant participation \cite{28_gonzalez2022allying,40_lee2017algorithmic}. These perceptions are not merely attitudinal; they have real consequences for behavior, as reduced trust in automated systems can lead qualified candidates to avoid applying altogether, particularly when processes lack transparency or feedback mechanisms. The implications are especially concerning for individuals from underrepresented or lower socioeconomic backgrounds, who are disproportionately affected by biased systems and structural inequities \cite{26_giannakos2017understanding,46_metaxa2018gender}. For these groups, negative experiences with automated hiring can influence not only immediate job outcomes but also long-term decisions about remaining in the field.

Finally, the persistence of social and structural inequalities within ostensibly “objective” automated systems provides a strong motivation for this study. Prior work has shown that hiring success is still significantly influenced by \textbf{social connections such as referrals and factors like socioeconomic status}, suggesting that AEDTs do not eliminate bias but may instead reshape or even reinforce existing inequities \cite{16_agyare2026why}. In some cases, applicants with greater social capital can effectively circumvent automated filters, while others rely on less effective “egalitarian” strategies that do not yield comparable outcomes. This raises critical questions about whether automation truly democratizes hiring or simply introduces new forms of advantage and exclusion. Taken together, these concerns motivate a comprehensive investigation into how perceptions of fairness, strategy use, and structural factors interact in automated hiring, with the goal of informing more equitable, transparent, and human-centered system design.

Qualitative research interviewing job-seekers has revealed the importance of referrals in circumventing automation and securing job offers, and that they often attribute their hiring successes to knowing how to “play the hiring game” \cite{4_armstrong2023navigating,16_agyare2026why}. Students from different socioeconomic backgrounds have reported different experiences with hiring processes even when they have access to the same resources, with applicants from working and middle-class backgrounds feeling like they need to earn social credit before relying on social connections \cite{16_agyare2026why}. Consequently, automated hiring experiences and knowledge of strategies may not only impact applicants desire to stay in the computing field, but also their ability to receive job offers. Further work is needed to assess applicants' experiences and perceptions of fairness in relation to socioeconomic status and resource access in the context of automated hiring processes. We expand on prior qualitative work by providing a quantitative survey investigation, examining how referrals and other strategies impact people's perceptions of procedural fairness, desire to be evaluated by AEDTs, and job-seeking success.

\section{Methodology}
\subsection{Participants and Recruitment}
The study recruited participants through a structured, multi-phase process targeting computer science students who are actively engaged in the job market. An initial power analysis indicated that a sample size between 170 and 250 participants would be sufficient to detect small to moderate effects, but the researchers exceeded this threshold by surveying a total of 525 undergraduate and graduate students across phases. The first phase focused on students from the University of Pennsylvania, where 294 participants were recruited from upper-level courses. The recruitment messaging clearly framed the study’s purpose as understanding fairness in automated hiring and emphasized its potential contribution to improving hiring practices and career support systems. Participants were informed about the nature of the survey, including its approximate duration (15 minutes) and the types of questions they would encounter, such as scenario-based evaluations of hiring processes and reflections on their own job search experiences.

To ensure ethical rigor and research validity, the study obtained approval from the Institutional Review Board (IRB) at the University of Pennsylvania, including approval for expanded recruitment in subsequent phases. Additionally, the researchers pre-registered their hypotheses and analysis plans on the Open Science Framework, reinforcing transparency and reproducibility in their methodology. The recruitment strategy aimed to capture a relevant population—young, technically trained job seekers—whose experiences are directly shaped by automated hiring systems. By combining careful sample size planning, multi-phase recruitment, and adherence to ethical and open science practices, the study establishes a strong methodological foundation for analyzing perceptions and strategies in automated hiring contexts.

\subsection{Creating Perceptions of Fairness}
The first part of the survey assessed participants' perceptions of procedural fairness in and willingness to be evaluated by 15 different hiring scenarios. These were made up of three evaluation types (technical coding assessment, resume reviewing, and behavioral interviews) and five different automation levels (human-only reviewing, human-AI reviewing equally, human reviewers for AI rejections, human reviewers for AI acceptances, AI-only reviewing).
\begin{enumerate}
    \item \textbf{Human-only reviewers}: An applicant submits a sample of code which is reviewed by a recruitment team who determines whether the applicant advances to the next phase.
    \item \textbf{Human-AI reviewing equally}: An applicant is given an online coding assessment, which is evaluated by an algorithm. If the applicant reaches a certain score on the autograder, the applicant advances to the next phase. All algorithmic decisions are reviewed by a recruitment team.
    \item \textbf{Human reviewers for AI rejections}: An applicant is given an online coding assessment, which is evaluated by an algorithm. If the algorithm rejects the applicant, the decision is reviewed by a recruitment team.
    \item \textbf{Human reviewers for AI acceptances}: An applicant is given an online coding assessment, which is evaluated by an algorithm. If the algorithm advances the applicant, the decision is reviewed by a recruitment team.
    \item \textbf{AI-only reviewing}: An applicant is given an online coding assessment, which is evaluated by an algorithm that determines whether an applicant advances to the next phase.
\end{enumerate}
The evaluation types were chosen based on prior work interviewing computing students about the types of automation they encountered in the hiring process \cite{4_armstrong2023navigating} and on the most prevalent automated employment decision tools (AEDTs) for the screening of candidates.

\subsection{Awareness of Experiences with Automated Hiring}
This section of the survey focused on understanding participants' awareness and direct experiences with automated hiring systems. Participants were asked whether they had encountered or were familiar with specific types of AEDTs, including online coding assessments, automated resume screening tools, and automated video interviews. For each, they selected from options indicating whether they had experienced it, heard of it, neither, or were unsure. In addition, the survey captured participants' perceptions of transparency and feedback by asking them to rate statements about understanding how their data was used in hiring and whether they received feedback from automated systems, using a five-point Likert\footnote{A Likert scale is a psychometric survey tool developed by Rensis Likert in 1932 to measure attitudes, opinions, and behaviors. It uses a 5-, 7-, or 9-point ordered range—typically from "Strongly Disagree" to "Strongly Agree" to quantify the intensity of a respondent's feelings, rather than just a binary yes/no.} scale.

The survey also examined the strategies participants used to navigate automated hiring processes. These included actions such as optimizing resumes with keywords, practicing coding problems using platforms like LeetCode, obtaining referrals, and leveraging personal or family connections in the industry. Participants further reported whether they were actively applying for jobs, the number of applications submitted, and whether they had received job offers; those not actively applying were excluded from analysis related to outcomes. Finally, optional demographic information—such as gender, race, and household income—was collected to explore how these factors might relate to job-seeking experiences and success.

\section{Key Findings}
\subsection{Participants’ Perceptions of the Hiring Process}
This section examines how participants perceive the fairness of automated hiring processes and their willingness to be evaluated by them, across different combinations of evaluation types and levels of automation. Participants rated both procedural fairness and willingness using Likert scales, and the results showed that these two measures closely aligned, although willingness was consistently slightly lower than perceived fairness. Overall, participants expressed skepticism toward automated hiring, indicating that while they may recognize some level of fairness, they are less personally comfortable being evaluated by such systems.

\subsubsection{Evaluation Type}
Participants' perceptions varied significantly depending on the type of evaluation being automated. Automation was viewed as more acceptable and fair in highly technical contexts, such as coding assessments, where evaluation criteria are more objective and measurable. In contrast, less technical and more subjective processes—particularly behavioral interviews and resume evaluations—were perceived as significantly less fair when automated. Participants emphasized that these processes involve nuance, context, and human judgment that algorithms struggle to capture, leading to concerns about misrepresentation of their abilities. Overall, there was a clear gradient: fairness and willingness were highest for coding assessments and lowest for interviews, highlighting that the appropriateness of automation depends heavily on the nature of the task.

\subsubsection{Automation Level}

The level of automation also played a critical role in shaping perceptions. Participants consistently preferred hiring processes that included some degree of human involvement, with fully automated (AI-only) systems being rated as the least fair and least desirable. Even when automation was present, scenarios where humans reviewed decisions—especially rejections—were perceived more favorably, suggesting that human oversight provides a sense of accountability and fairness. As the level of automation increased and human involvement decreased, both fairness perceptions and willingness to participate declined. This indicates that participants do not reject automation entirely but are strongly opposed to fully automated decision-making without human intervention.

\subsection{Strategy Use}
This section examines how applicants use different strategies to navigate automated hiring systems and how these strategies relate to both their perceptions of fairness and their job outcomes. The analysis connects strategy use with two key research questions: whether strategic behavior influences perceptions of AEDTs, and whether awareness and strategy use translate into better hiring success. Overall, the findings suggest that while applicants actively engage in multiple strategies, these do not uniformly improve perceptions or outcomes, highlighting a disconnect between effort and effectiveness in automated hiring environments.

\subsubsection{Descriptive Results}
Participants reported widespread use of various strategies to navigate automated hiring processes, with some strategies being significantly more common than others. The most frequently used approaches included leveraging online resources, practicing coding problems (e.g., via LeetCode), and consulting peers who had recently gone through the hiring process. Many participants also engaged in mass-applying to jobs and modifying resumes with keywords, although fewer used resume scanners to test how their applications would be parsed. Referrals were also common, with over half of participants receiving at least one, though the extent of reliance varied widely. Additionally, participants demonstrated high awareness and exposure to AEDTs, especially coding assessments and video interviews, but reported limited understanding of how their data was used and a lack of feedback from these systems.

\subsubsection{Strategy Use and Participant Perceptions}
The relationship between strategy use and perceptions of fairness revealed a somewhat counterintuitive pattern. Participants who used certain strategies—particularly referrals or who were more aware of coding assessment tools—tended to report lower perceptions of fairness and lower willingness to be evaluated by AEDTs. This suggests that greater familiarity with the hiring process may expose its limitations or “game-like” nature, leading to increased skepticism. In contrast, awareness of automated video interviews was associated with slightly higher fairness and willingness ratings. Notably, many commonly used strategies—such as resume keyword optimization, practicing coding problems, mass-applying, or using career services—did not have a significant impact on fairness perceptions. Demographic factors also did not significantly influence these perceptions, indicating that strategy-related awareness, rather than background characteristics, played a larger role in shaping attitudes.

\subsubsection{Strategy Use and Job Outcomes}
When examining the impact of strategies on actual hiring success, the findings showed that most commonly used strategies did not significantly improve outcomes. The only strategy that demonstrated a meaningful positive relationship with receiving job offers was the use of referrals, highlighting the continued importance of social connections in hiring. Additionally, family household income emerged as another strong predictor of success, reinforcing the role of socioeconomic factors in shaping job market outcomes. In contrast, widely adopted “egalitarian” strategies—such as practicing coding problems or optimizing resumes—did not translate into increased job success. These results suggest that while applicants invest effort into navigating automated systems, structural advantages like social capital play a far more decisive role in determining outcomes.
